%!TEX program = pdflatex
\documentclass[11pt,twoside,openright]{book}

% ─── Packages ──────────────────────────────────────────────────────
\usepackage[margin=1in,inner=1.3in,outer=1in]{geometry}
\usepackage{amsmath,amssymb,amsthm,mathtools}
\usepackage{tikz}
\usetikzlibrary{calc,arrows.meta,decorations.markings,patterns,
  shapes.geometric,positioning,3d,perspective,math,shadows,
  decorations.pathmorphing,backgrounds,fit}
\usepackage{pgfplots}
\pgfplotsset{compat=1.18}
\usepackage{graphicx}
\usepackage{xcolor}
\usepackage{hyperref}
\usepackage{enumitem}
\usepackage{fancyhdr}
\usepackage{titlesec}
\usepackage{epigraph}
\usepackage{listings}
\usepackage{microtype}
\usepackage{booktabs}
\usepackage{multirow}
\usepackage{tcolorbox}
\tcbuselibrary{listings,skins,breakable}
\usepackage{fontenc}
\usepackage[utf8]{inputenc}
\usepackage{lettrine}
\usepackage{setspace}

% ─── Colors ────────────────────────────────────────────────────────
\definecolor{deepblue}{RGB}{0,51,102}
\definecolor{goldaccent}{RGB}{204,153,0}
\definecolor{darkgreen}{RGB}{0,100,0}
\definecolor{codebg}{RGB}{248,248,242}
\definecolor{codeframe}{RGB}{180,180,180}
\definecolor{leanblue}{RGB}{30,80,140}
\definecolor{pythagorean}{RGB}{139,69,19}
\definecolor{lorentzred}{RGB}{178,34,34}
\definecolor{tropicalgreen}{RGB}{34,139,34}
\definecolor{fermatpurple}{RGB}{102,0,153}
\definecolor{cayleyorange}{RGB}{204,85,0}
\definecolor{chapterbg}{RGB}{240,240,255}

% ─── Theorem environments ─────────────────────────────────────────
\theoremstyle{definition}
\newtheorem{theorem}{Theorem}[chapter]
\newtheorem{lemma}[theorem]{Lemma}
\newtheorem{proposition}[theorem]{Proposition}
\newtheorem{corollary}[theorem]{Corollary}
\newtheorem{definition}[theorem]{Definition}
\newtheorem{example}[theorem]{Example}
\newtheorem{remark}[theorem]{Remark}

% ─── Lean code listing style ──────────────────────────────────────
\lstdefinelanguage{Lean}{
  keywords={theorem,lemma,def,structure,class,instance,where,
    import,open,section,namespace,end,variable,
    by,have,let,show,calc,exact,apply,intro,intros,
    rfl,ring,omega,simp,norm_num,linarith,nlinarith,
    sorry,example,noncomputable,if,then,else,match,with,
    fun,return,do,for,in,inductive,axiom},
  sensitive=true,
  morecomment=[l]{--},
  morecomment=[s]{/-}{-/},
  morestring=[b]",
  literate={→}{{$\to$}}1 {←}{{$\leftarrow$}}1 {↦}{{$\mapsto$}}1
           {∀}{{$\forall$}}1 {∃}{{$\exists$}}1 {¬}{{$\neg$}}1
           {∧}{{$\wedge$}}1 {∨}{{$\lor$}}1 {≠}{{$\neq$}}1
           {≤}{{$\leq$}}1 {≥}{{$\geq$}}1 {⊤}{{$\top$}}1
           {ℤ}{{$\mathbb{Z}$}}1 {ℕ}{{$\mathbb{N}$}}1
           {ℝ}{{$\mathbb{R}$}}1 {ℂ}{{$\mathbb{C}$}}1
           {ℍ}{{$\mathbb{H}$}}1 {ℚ}{{$\mathbb{Q}$}}1
           {α}{{$\alpha$}}1 {β}{{$\beta$}}1 {γ}{{$\gamma$}}1
           {ε}{{$\varepsilon$}}1 {ζ}{{$\zeta$}}1 {ω}{{$\omega$}}1
           {λ}{{$\lambda$}}1 {σ}{{$\sigma$}}1
           {×}{{$\times$}}1 {⟨}{{$\langle$}}1 {⟩}{{$\rangle$}}1
           {∣}{{$\mid$}}1 {·}{{$\cdot$}}1
           {ᵀ}{{${}^{\top}$}}1 {₁}{{$_1$}}1 {₂}{{$_2$}}1
           {₃}{{$_3$}}1 {₄}{{$_4$}}1 {₀}{{$_0$}}1
           {⁻}{{${}^{-}$}}1 {²}{{${}^2$}}1 {³}{{${}^3$}}1
           {⁴}{{${}^4$}}1 {π}{{$\pi$}}1 {μ}{{$\mu$}}1
           {↔}{{$\leftrightarrow$}}1 {⊢}{{$\vdash$}}1
           {⟩}{{$\rangle$}}1 {⟨}{{$\langle$}}1
           {▸}{{$\blacktriangleright$}}1
           {⊥}{{$\bot$}}1 {♯}{{\#}}1 {√}{{$\sqrt{}$}}1
           {≡}{{$\equiv$}}1 {⊔}{{$\sqcup$}}1
           {⊓}{{$\sqcap$}}1 {∩}{{$\cap$}}1 {∪}{{$\cup$}}1
           {∈}{{$\in$}}1 {∉}{{$\notin$}}1 {⊆}{{$\subseteq$}}1
           {∅}{{$\emptyset$}}1 {∑}{{$\sum$}}1 {∏}{{$\prod$}}1
           {↑}{{$\uparrow$}}1 {↓}{{$\downarrow$}}1
           {⌊}{{$\lfloor$}}1 {⌋}{{$\rfloor$}}1
           {⌈}{{$\lceil$}}1 {⌉}{{$\rceil$}}1
           {∘}{{$\circ$}}1 {𝕆}{{$\mathbb{O}$}}1
           {𝕊}{{$\mathbb{S}$}}1,
}

\lstset{
  language=Lean,
  basicstyle=\ttfamily\footnotesize,
  keywordstyle=\color{leanblue}\bfseries,
  commentstyle=\color{darkgreen}\itshape,
  stringstyle=\color{pythagorean},
  backgroundcolor=\color{codebg},
  frame=single,
  rulecolor=\color{codeframe},
  breaklines=true,
  breakatwhitespace=false,
  tabsize=2,
  showstringspaces=false,
  numbers=left,
  numberstyle=\tiny\color{gray},
  numbersep=5pt,
  xleftmargin=1.5em,
  framexleftmargin=1em,
  captionpos=b,
  extendedchars=true,
  inputencoding=utf8,
}

% ─── tcolorbox styles ─────────────────────────────────────────────
\newtcolorbox{insightbox}[1][]{
  colback=chapterbg, colframe=deepblue,
  fonttitle=\bfseries, title={#1},
  arc=3mm, boxrule=0.8pt,
  breakable
}

\newtcolorbox{marginalia}[1][]{
  colback=yellow!5, colframe=goldaccent,
  fonttitle=\bfseries\itshape, title={#1},
  arc=2mm, boxrule=0.5pt,
  left=5pt, right=5pt, top=3pt, bottom=3pt
}

% ─── Chapter title formatting ─────────────────────────────────────
\titleformat{\chapter}[display]
  {\normalfont\huge\bfseries\color{deepblue}}
  {\chaptertitlename\ \thechapter}{20pt}{\Huge}
\titlespacing*{\chapter}{0pt}{-30pt}{30pt}

% ─── Headers ──────────────────────────────────────────────────────
\pagestyle{fancy}
\fancyhf{}
\fancyhead[LE]{\small\textit{\leftmark}}
\fancyhead[RO]{\small\textit{\rightmark}}
\fancyfoot[C]{\thepage}
\renewcommand{\headrulewidth}{0.4pt}

% ─── Epigraph style ──────────────────────────────────────────────
\setlength{\epigraphwidth}{0.7\textwidth}
\renewcommand{\epigraphflush}{center}
\renewcommand{\epigraphrule}{0pt}
\renewcommand{\textflush}{flushepinormal}

% ─── Hyperref ─────────────────────────────────────────────────────
\hypersetup{
  colorlinks=true,
  linkcolor=deepblue,
  urlcolor=deepblue,
  citecolor=darkgreen,
  pdftitle={The Eternal Braid: Pythagoras, Lorentz, and the Architecture of Proof},
  pdfauthor={Paul Klemstine},
}

% ─── Custom commands ──────────────────────────────────────────────
\newcommand{\Q}{\mathbb{Q}}
\newcommand{\Z}{\mathbb{Z}}
\newcommand{\R}{\mathbb{R}}
\newcommand{\N}{\mathbb{N}}
\newcommand{\C}{\mathbb{C}}
\newcommand{\HH}{\mathbb{H}}
\newcommand{\OO}{\mathbb{O}}
\newcommand{\Fin}{\mathrm{Fin}}
\newcommand{\lorentzQ}{\mathcal{Q}}
\newcommand{\sorry}{\texttt{sorry}}

% ═══════════════════════════════════════════════════════════════════
\begin{document}

% ─── Half-title page ──────────────────────────────────────────────
\begin{titlepage}
\thispagestyle{empty}
\vspace*{3cm}
\begin{center}
{\Huge\bfseries\color{deepblue} The Eternal Braid}\\[1.5cm]
{\Large\itshape Pythagoras, Lorentz, and the\\[0.3cm]
Architecture of Proof}
\end{center}
\vfill
\end{titlepage}

\cleardoublepage

% ─── Full title page ─────────────────────────────────────────────
\begin{titlepage}
\thispagestyle{empty}
\vspace*{1cm}

\begin{center}

% ─── Decorative header ───
\begin{tikzpicture}
  % Golden spiral approximation
  \draw[goldaccent, line width=1.5pt]
    (0,0) -- (1,0) arc(0:90:1) -- (-1,1) arc(90:180:1)
    -- (-2,-1) arc(180:270:2) -- (1,-3) arc(270:360:3);
  % Three interlocking circles (Borromean-inspired)
  \draw[deepblue, line width=1pt] (0,0) circle(0.8);
  \draw[lorentzred, line width=1pt] (0.7,-0.4) circle(0.8);
  \draw[darkgreen, line width=1pt] (0.35,0.6) circle(0.8);
\end{tikzpicture}

\vspace{1cm}

{\fontsize{36}{42}\selectfont\bfseries\color{deepblue}
The Eternal Braid}\\[0.8cm]

{\fontsize{18}{22}\selectfont\itshape\color{gray}
Pythagoras, Lorentz, and the\\[0.2cm]
Architecture of Proof}\\[1.5cm]

\rule{8cm}{0.4pt}\\[1cm]

{\Large\scshape How Ancient Triangles Encode\\[0.2cm]
Relativistic Geometry, Quantum Shortcuts,\\[0.2cm]
and the Secrets of Prime Numbers}\\[1cm]

\rule{8cm}{0.4pt}\\[1.5cm]

{\LARGE Paul Klemstine}\\[0.3cm]
{\large\itshape Independent Researcher in Mathematics and AI}\\[2cm]

{\large With Machine-Verified Proofs in Lean 4}\\[0.5cm]
{\small\itshape All theorems formally checked by computer --- no margin too small}

\vfill

{\small\itshape First Edition, 2026}

\end{center}
\end{titlepage}

\cleardoublepage

% ─── Dedication ───────────────────────────────────────────────────
\thispagestyle{empty}
\vspace*{5cm}
\begin{center}
{\Large\itshape Soli Deo Gloria}\\[1cm]
{\normalsize To God alone be the glory.}\\[2cm]
{\small\itshape ``The heavens declare the glory of God;\\
the skies proclaim the work of his hands.\\
Day after day they pour forth speech;\\
night after night they reveal knowledge.''}\\[0.3cm]
{\small --- Psalm 19:1--2}
\end{center}
\vfill
\cleardoublepage

% ─── Epigraph page ───────────────────────────────────────────────
\thispagestyle{empty}
\vspace*{4cm}
\begin{center}
{\itshape ``In mathematics you don't understand things.\\
You just get used to them.''}\\[0.3cm]
{\small --- John von Neumann}\\[2cm]

{\itshape ``A mathematician, like a painter or a poet, is a maker of patterns.\\
If his patterns are more permanent than theirs, it is because\\
they are made with ideas.''}\\[0.3cm]
{\small --- G. H. Hardy, \emph{A Mathematician's Apology}}\\[2cm]

{\itshape ``The question you raise, `how can such a formulation lead to\\
computations?' doesn't bother me in the least.\\
Throughout my whole life as a mathematician, the possibility\\
of making explicit, elegant computations has always come\\
out by itself, as a byproduct of a thorough conceptual\\
understanding of what was going on.''}\\[0.3cm]
{\small --- Grothendieck}
\end{center}
\vfill
\cleardoublepage

% ─── Table of Contents ───────────────────────────────────────────
\tableofcontents
\cleardoublepage

% ═══════════════════════════════════════════════════════════════════
%  PRELUDE
% ═══════════════════════════════════════════════════════════════════
\chapter*{Prelude: The Margin and the Machine}
\addcontentsline{toc}{chapter}{Prelude: The Margin and the Machine}
\markboth{Prelude}{The Margin and the Machine}

\lettrine[lines=3,loversize=0.15]{I}{n 1637}, Pierre de Fermat scribbled
a note in the margin of his copy of Diophantus's \emph{Arithmetica}. He had
discovered, he claimed, ``a truly marvelous proof'' that no three positive
integers $a$, $b$, $c$ can satisfy $a^n + b^n = c^n$ for any integer $n > 2$.
The margin, famously, was too narrow to contain it.

Three and a half centuries later, Andrew Wiles filled 109 pages---not of a
margin, but of the \emph{Annals of Mathematics}---to prove Fermat right.
The proof required modular forms, Galois representations, and the full
machinery of 20th-century algebraic geometry. No margin could ever have held it.

But this book is not about Fermat's Last Theorem, not really. It is about a
different kind of margin: the gap between what humans can verify and what
machines can prove. It is about an ancient equation---$a^2 + b^2 = c^2$---and
the astonishing web of connections that radiate outward from it, touching
relativity, quantum mechanics, cryptography, and the very foundations of
mathematical truth.

Every theorem in this book has been \emph{machine-verified} using Lean 4,
a proof assistant developed at Microsoft Research. This means that a computer
has checked every logical step, every case analysis, every algebraic
manipulation. The margin is no longer too small: the margin is infinite,
and it is inhabited by silicon.

\begin{center}
\begin{tikzpicture}[scale=0.9]
  % Three pillars
  \fill[deepblue!10] (-4,-0.5) rectangle (-2,4);
  \fill[lorentzred!10] (-1,-0.5) rectangle (1,4);
  \fill[darkgreen!10] (2,-0.5) rectangle (4,4);

  \node[font=\small\bfseries,color=deepblue] at (-3,4.3) {Pythagoras};
  \node[font=\small\bfseries,color=lorentzred] at (0,4.3) {Lorentz};
  \node[font=\small\bfseries,color=darkgreen] at (3,4.3) {Lean};

  \node[font=\footnotesize,text width=1.8cm,align=center] at (-3,3) {$a^2+b^2$\\$=c^2$};
  \node[font=\footnotesize,text width=1.8cm,align=center] at (0,3) {$ds^2=$\\$dt^2-dx^2$};
  \node[font=\footnotesize,text width=1.8cm,align=center] at (3,3) {\texttt{theorem}\\proved by\\kernel};

  % Connecting arcs
  \draw[->,thick,goldaccent] (-2,2) to[bend left=20] (-1,2);
  \draw[->,thick,goldaccent] (1,2) to[bend left=20] (2,2);
  \draw[->,thick,goldaccent] (-2,1) to[bend right=30] (2,1);

  \node[font=\tiny,goldaccent] at (-1.5,2.5) {isometry};
  \node[font=\tiny,goldaccent] at (1.5,2.5) {verification};
  \node[font=\tiny,goldaccent] at (0,0.5) {the eternal braid};
\end{tikzpicture}
\end{center}

\medskip

The book you hold braids three strands together:

\begin{enumerate}[leftmargin=2em]
\item \textbf{The Pythagorean Strand.} The equation $a^2+b^2=c^2$ generates
an infinite ternary tree---the Berggren tree---whose structure encodes deep
secrets about prime numbers and integer factoring.

\item \textbf{The Lorentz Strand.} The matrices that generate this tree are
elements of $O(2,1;\Z)$, the integer Lorentz group from special relativity.
The tree tiles the hyperbolic plane. Pythagorean triple generation is,
secretly, relativistic geometry.

\item \textbf{The Proof Strand.} Every major theorem is formalized in Lean 4
with Mathlib, creating an unbreakable chain of logical certainty from axioms
to conclusions. The machine has checked the math so you don't have to---but
you might want to anyway, because the math is beautiful.
\end{enumerate}

Like Hofstadter's \emph{G\"odel, Escher, Bach}, this book is about
\emph{strange loops}---patterns that, when followed far enough, fold back
on themselves in unexpected ways. A right triangle becomes a point in
hyperbolic space. A matrix multiplication becomes a walk through a tree.
A tree traversal becomes a factoring algorithm. A factoring algorithm
becomes a question about the security of the internet. And a question
about security becomes, once again, a question about right triangles.

The braid is eternal. Let us begin to trace its strands.

% ═══════════════════════════════════════════════════════════════════
%  PART I: THE TREE
% ═══════════════════════════════════════════════════════════════════
\part{The Tree of Triples}

% ───────────────────────────────────────────────────────────────────
\chapter{The Berggren--Lorentz Correspondence}
\label{ch:berggren-lorentz}

\epigraph{\itshape ``There is geometry in the humming of the strings,\\
there is music in the spacing of the spheres.''}
{--- attributed to Pythagoras}

\lettrine[lines=3]{E}{very schoolchild} knows the triple $(3,4,5)$: three
squared plus four squared equals five squared, and you can build a right
triangle from a piece of rope with 12 equally-spaced knots. What fewer
people know is that $(3,4,5)$ is the \emph{root} of an infinite tree that
contains \emph{every} primitive Pythagorean triple exactly once.

This tree was discovered by Berggren in 1934, and it harbors a
secret that takes us straight from ancient geometry into the heart of
Einstein's special relativity.

\section{The Null Cone}

Consider the quadratic form
\[
\lorentzQ(a,b,c) = a^2 + b^2 - c^2.
\]
In physics, this is the \emph{Lorentz form}---it measures the ``interval''
between events in spacetime (with signature $++-$ instead of the usual
$+--$). A Pythagorean triple $(a,b,c)$ satisfies $a^2+b^2=c^2$, which
is precisely the condition $\lorentzQ(a,b,c) = 0$.

\begin{insightbox}[The Null Cone Principle]
Pythagorean triples are the \emph{integer points on the null cone}
of the Lorentz form. They are the lattice points where a light ray
passes through, in a two-plus-one-dimensional Minkowski spacetime
with integer coordinates.
\end{insightbox}

\begin{center}
\begin{tikzpicture}[scale=1.1]
  % Draw the null cone
  \begin{scope}[canvas is xz plane at y=0]
    \draw[lorentzred!30, fill=lorentzred!5] (0,0) -- (3,3) arc(0:360:3 and 3) -- cycle;
  \end{scope}

  % Cone outline
  \draw[lorentzred, thick] (0,0) -- (2.5,3.5);
  \draw[lorentzred, thick] (0,0) -- (-2.5,3.5);
  \draw[lorentzred, thick, dashed] (0,0) -- (2.5,-0.5);
  \draw[lorentzred, thick, dashed] (0,0) -- (-2.5,-0.5);

  % Ellipse at top
  \draw[lorentzred, thick] (0,3.5) ellipse(2.5 and 0.6);

  % Lattice points on cone
  \fill[deepblue] (0.9,1.5) circle(2.5pt) node[right,font=\tiny] {$(3,4,5)$};
  \fill[deepblue] (-0.6,2.0) circle(2.5pt) node[left,font=\tiny] {$(5,12,13)$};
  \fill[deepblue] (1.5,2.5) circle(2.5pt) node[right,font=\tiny] {$(8,15,17)$};
  \fill[deepblue] (-1.2,2.8) circle(2.5pt) node[left,font=\tiny] {$(7,24,25)$};
  \fill[deepblue] (0.3,3.0) circle(2.5pt) node[right,font=\tiny] {$(20,21,29)$};

  % Labels
  \node[lorentzred, font=\small\bfseries] at (3.2,3.5) {$\lorentzQ = 0$};
  \node[font=\small] at (0,-1.2) {The Null Cone of $\lorentzQ(a,b,c) = a^2+b^2-c^2$};

  % Axes
  \draw[->,thick] (0,0) -- (3.5,0) node[right,font=\small] {$a$};
  \draw[->,thick] (0,0) -- (0,4.2) node[above,font=\small] {$c$};
\end{tikzpicture}
\end{center}

\section{Three Matrices, One Group}

Berggren discovered that three $3\times 3$ integer matrices generate every
primitive Pythagorean triple from the root $(3,4,5)$:
\[
B_A = \begin{pmatrix} 1 & -2 & 2 \\ 2 & -1 & 2 \\ 2 & -2 & 3\end{pmatrix},\quad
B_B = \begin{pmatrix} 1 & 2 & 2 \\ 2 & 1 & 2 \\ 2 & 2 & 3\end{pmatrix},\quad
B_C = \begin{pmatrix} -1 & 2 & 2 \\ -2 & 1 & 2 \\ -2 & 2 & 3\end{pmatrix}.
\]

The astonishing fact, which we prove with machine-verified certainty, is:
\[
B_X^{\!\top}\, Q\, B_X = Q, \qquad X \in \{A, B, C\},
\]
where $Q = \mathrm{diag}(1,1,-1)$ is the Lorentz metric. These matrices
\emph{preserve the Lorentz form}. They are elements of the integer Lorentz
group $O(2,1;\Z)$.

\begin{center}
\begin{tikzpicture}[
  level distance=1.8cm,
  sibling distance=3.5cm,
  edge from parent/.style={draw,thick,->,deepblue},
  every node/.style={draw,rounded corners,fill=deepblue!8,
    font=\small,minimum width=1.6cm,minimum height=0.6cm}
]
  \node {$3,4,5$}
    child {node {$5,12,13$}
      child {node {$7,24,25$} edge from parent node[left,draw=none,fill=none,font=\tiny] {$A$}}
      child {node {$45,28,53$} edge from parent node[left,draw=none,fill=none,font=\tiny] {$B$}}
      child {node {$55,48,73$} edge from parent node[right,draw=none,fill=none,font=\tiny] {$C$}}
      edge from parent node[left,draw=none,fill=none,font=\tiny] {$A$}
    }
    child {node {$21,20,29$}
      child {node[font=\tiny] {$\cdots$} edge from parent}
      child {node[font=\tiny] {$\cdots$} edge from parent}
      child {node[font=\tiny] {$\cdots$} edge from parent}
      edge from parent node[right,draw=none,fill=none,font=\tiny\color{black}] {$B$}
    }
    child {node {$15,8,17$}
      child {node[font=\tiny] {$\cdots$} edge from parent}
      child {node[font=\tiny] {$\cdots$} edge from parent}
      child {node[font=\tiny] {$\cdots$} edge from parent}
      edge from parent node[right,draw=none,fill=none,font=\tiny\color{black}] {$C$}
    };
\end{tikzpicture}
\end{center}

\section{The Factoring Bridge}

Here is where the ancient becomes urgent. For any Pythagorean triple
$(a,b,c)$ with $a^2 + b^2 = c^2$:
\[
(c-b)(c+b) = a^2.
\]
This is the \emph{difference of squares identity}, and it is the
foundation of every modern factoring algorithm. If $a^2 = N$ and we
can find $b,c$ with $b^2 \equiv c^2 \pmod{N}$, then $\gcd(c-b, N)$
often reveals a factor.

The Berggren tree provides a systematic way to search for such
representations, transforming factoring into tree traversal.

\begin{insightbox}[The Pell Recurrence]
The $B$-branch of the Berggren tree satisfies a three-term recurrence
on hypotenuses: $c_{n+2} = 6c_{n+1} - c_n$. This is a \emph{Pell
equation} recurrence, connecting Pythagorean triples to the deepest
structures of algebraic number theory.
\end{insightbox}

\section{What the Machine Says}

In Lean 4, the Lorentz preservation theorem is stated and proved as:

\begin{lstlisting}[caption={Berggren--Lorentz preservation (machine-verified)}]
theorem berggrenA_lorentz :
    berggrenA_matrix ^ * lorentzMetric * berggrenA_matrix
    = lorentzMetric := by
  native_decide
\end{lstlisting}

The keyword \texttt{native\_decide} instructs Lean to compile the
computation to native code and verify the matrix equation by brute force.
The kernel checks every entry. There is no room for error.

% ───────────────────────────────────────────────────────────────────
\chapter{The Lattice--Tree Correspondence}
\label{ch:lattice-tree}

\epigraph{\itshape ``God made the integers; all else is the work of man.''}
{--- Leopold Kronecker}

\lettrine[lines=3]{T}{he most} surprising theorem in this book is perhaps the
simplest to state: \emph{walking down the Berggren tree is the same algorithm
as running the Euclidean algorithm on the parameters of a Pythagorean triple.}

\section{Berggren Descent as Lattice Reduction}

Every primitive Pythagorean triple can be written as
$(m^2-n^2,\; 2mn,\; m^2+n^2)$ for coprime $m > n > 0$ of opposite parity
(the Euclid parametrization). The Berggren tree inverse matrices act on
the parameter pair $(m,n)$ via $2\times 2$ matrices in $\mathrm{SL}(2,\Z)$:
\[
M_1 = \begin{pmatrix} 2 & -1 \\ 1 & 0 \end{pmatrix},\quad
M_3 = \begin{pmatrix} 1 & 2 \\ 0 & 1 \end{pmatrix}.
\]

The inverse $M_3^{-1}$ maps $(m,n) \mapsto (m - 2n, n)$---exactly the
subtraction step of the Euclidean algorithm with quotient $q = 2$. The
inverse $M_1^{-1}$ performs the swap: $(m,n) \mapsto (n, 2n-m)$.

\begin{center}
\begin{tikzpicture}[scale=1.0]
  % Euclidean algorithm flow
  \node[draw, rounded corners, fill=deepblue!10, minimum width=2.8cm] (a) at (0,4) {$(m, n)$};
  \node[draw, rounded corners, fill=deepblue!10, minimum width=2.8cm] (b) at (0,2.5) {$(m{-}2n, n)$};
  \node[draw, rounded corners, fill=goldaccent!20, minimum width=2.8cm] (c) at (0,1) {$(n, 2n{-}m)$};
  \node[draw, rounded corners, fill=deepblue!10, minimum width=2.8cm] (d) at (0,-0.5) {$\vdots$};

  \draw[->,thick] (a) -- (b) node[midway, right, font=\small] {$M_3^{-1}$: subtract};
  \draw[->,thick] (b) -- (c) node[midway, right, font=\small] {$M_1^{-1}$: swap};
  \draw[->,thick] (c) -- (d) node[midway, right, font=\small] {repeat};

  % Berggren tree flow
  \node[draw, rounded corners, fill=lorentzred!10, minimum width=2.8cm] (e) at (6,4) {$(a,b,c)$};
  \node[draw, rounded corners, fill=lorentzred!10, minimum width=2.8cm] (f) at (6,2.5) {$(a',b',c')$};
  \node[draw, rounded corners, fill=lorentzred!10, minimum width=2.8cm] (g) at (6,1) {$(a'',b'',c'')$};
  \node[draw, rounded corners, fill=lorentzred!10, minimum width=2.8cm] (h) at (6,-0.5) {$\vdots$};

  \draw[->,thick] (e) -- (f) node[midway, right, font=\small] {$B_X^{-1}$};
  \draw[->,thick] (f) -- (g) node[midway, right, font=\small] {$B_Y^{-1}$};
  \draw[->,thick] (g) -- (h) node[midway, right, font=\small] {descend};

  % Equivalence arrows
  \draw[<->,very thick, goldaccent, dashed] (2.3,3) -- (3.7,3)
    node[midway, above, font=\small\bfseries\color{goldaccent}] {$\equiv$};
  \draw[<->,very thick, goldaccent, dashed] (2.3,1.5) -- (3.7,1.5);

  % Labels
  \node[font=\bfseries, deepblue] at (0,5) {Euclidean Algorithm};
  \node[font=\bfseries, lorentzred] at (6,5) {Berggren Descent};
\end{tikzpicture}
\end{center}

\section{The Complexity Consequence}

Since the Euclidean algorithm runs in $O(\log \max(m,n))$ steps, and
for a balanced semiprime $N = pq$ we have $m \sim \sqrt{N}$, the
Berggren tree descent takes $\Theta(\sqrt{N})$ steps. This is
\emph{exactly} the complexity of trial division.

\begin{theorem}[Lattice--Tree Correspondence]
Inverse Berggren tree traversal computes exactly the same sequence
of quotients as the continued fraction expansion of $m/n$.
Consequently, Pythagorean tree factoring requires $\Theta(\sqrt{N})$
operations for balanced semiprimes $N = pq$.
\end{theorem}

This is both a disappointment and a revelation. The disappointment:
the Berggren tree, for all its beauty, cannot beat trial division in
two dimensions. The revelation: the proof identifies the \emph{escape
route}---higher-dimensional lattices, where LLL reduction can find
short vectors that Gauss-like methods miss.

% ───────────────────────────────────────────────────────────────────
\chapter{Hyperbolic Shortcuts}
\label{ch:hyperbolic}

\epigraph{\itshape ``Not all those who wander are lost.''}
{--- J.R.R. Tolkien}

\lettrine[lines=3]{T}{he Berggren tree} lives in hyperbolic space, and
hyperbolic space has an extraordinary property: you can take
\emph{shortcuts}. Path concatenation in the tree corresponds to matrix
multiplication, and matrix multiplication can be done in $O(\log k)$
steps by repeated squaring.

\section{Paths as Matrix Products}

Define a \emph{path} in the Berggren tree as a finite sequence of
branch labels $A$, $B$, $C$. The matrix corresponding to a path
$p_1 p_2 \cdots p_k$ is:
\[
M(p_1 p_2 \cdots p_k) = B_{p_1}\, B_{p_2}\, \cdots\, B_{p_k}.
\]

The key compositionality theorem states:
\[
M(p \cdot q) = M(p) \cdot M(q),
\]
where $p \cdot q$ is path concatenation. This means we can jump to any
node at depth $k$ in $O(\log k)$ matrix multiplications.

\begin{center}
\begin{tikzpicture}[scale=0.85]
  % Hyperbolic disk (Poincar\'e model)
  \fill[blue!3] (0,0) circle(3.5);
  \draw[deepblue, thick] (0,0) circle(3.5);

  % Geodesics (arcs in the Poincar\'e disk)
  \draw[deepblue!40, thick] (-3.5,0) arc(180:0:3.5 and 5);
  \draw[lorentzred!40, thick] (0,-3.5) arc(270:90:5 and 3.5);
  \draw[darkgreen!40, thick] (-2.47,-2.47) arc(225:45:3.5);

  % Tree nodes
  \fill[deepblue] (0,0) circle(3pt) node[above right,font=\small] {root};
  \fill[lorentzred] (1.5,1.2) circle(2.5pt);
  \fill[lorentzred] (-1.0,1.8) circle(2.5pt);
  \fill[lorentzred] (0.5,-1.5) circle(2.5pt);
  \fill[darkgreen] (2.2,0.5) circle(2pt);
  \fill[darkgreen] (2.5,1.8) circle(2pt);
  \fill[darkgreen] (-2.0,1.0) circle(2pt);
  \fill[darkgreen] (-0.5,2.5) circle(2pt);
  \fill[darkgreen] (1.5,-2.0) circle(2pt);
  \fill[darkgreen] (-1.0,-2.2) circle(2pt);

  % Tree edges
  \draw[thick] (0,0) -- (1.5,1.2);
  \draw[thick] (0,0) -- (-1.0,1.8);
  \draw[thick] (0,0) -- (0.5,-1.5);
  \draw[thin] (1.5,1.2) -- (2.2,0.5);
  \draw[thin] (1.5,1.2) -- (2.5,1.8);
  \draw[thin] (-1.0,1.8) -- (-2.0,1.0);
  \draw[thin] (-1.0,1.8) -- (-0.5,2.5);
  \draw[thin] (0.5,-1.5) -- (1.5,-2.0);
  \draw[thin] (0.5,-1.5) -- (-1.0,-2.2);

  % Shortcut arrow
  \draw[->,very thick, goldaccent, dashed, decorate,
    decoration={snake,amplitude=1.5pt,segment length=8pt}]
    (0,0) to[bend right=20] (2.5,1.8);
  \node[goldaccent, font=\small\bfseries] at (2.0,-0.3) {shortcut!};

  \node[font=\small] at (0,-4.3) {The Berggren tree tiling the Poincar\'e disk};
\end{tikzpicture}
\end{center}

\section{Inverse Matrices for Ascent}

Each Berggren matrix has an integer inverse (since $\det(B_A) = \det(B_C) = 1$
and $\det(B_B) = -1$). This means we can \emph{ascend} the tree---given any
primitive Pythagorean triple, we can determine which branch of the tree it
lives on by checking which inverse matrix produces a valid (positive-coordinate)
triple.

\begin{theorem}[Deterministic Descent]
At most one inverse Berggren matrix produces a triple with all positive
coordinates. Tree descent is deterministic: there is a unique path from
any triple back to the root $(3,4,5)$.
\end{theorem}

\section{The Chebyshev Connection}

The $B$-branch powers satisfy a Chebyshev recurrence. The hypotenuse
values along repeated $B$-steps follow:
\[
c_{n+1} = 6c_n - c_{n-1},
\]
which is exactly the recurrence for Chebyshev polynomials evaluated at $x=3$.
This connects the Berggren tree to the theory of orthogonal polynomials
and, through them, to approximation theory and spectral methods.

% ───────────────────────────────────────────────────────────────────
\chapter{Three Roads from Pythagoras}
\label{ch:three-roads}

\epigraph{\itshape ``Do not worry about your difficulties in mathematics.\\
I can assure you mine are still greater.''}
{--- Albert Einstein}

\lettrine[lines=3]{F}{rom the} Pythagorean equation $a^2+b^2=c^2$, three
distinct roads lead to integer factoring. Each road uses a different
mathematical mechanism, and each illuminates a different facet of the
factoring problem.

\section{Road I: Euler's Method}

If a number $N$ can be written as a sum of two squares in \emph{two different ways}:
\[
N = a^2 + b^2 = c^2 + d^2,
\]
then we can extract a factor of $N$ as:
\[
\gcd(a \cdot c + b \cdot d,\; N).
\]

\begin{center}
\begin{tikzpicture}[scale=0.9]
  % Two representations
  \node[draw, fill=deepblue!10, rounded corners, minimum width=3cm]
    (rep1) at (-2.5,2) {$N = a^2 + b^2$};
  \node[draw, fill=lorentzred!10, rounded corners, minimum width=3cm]
    (rep2) at (2.5,2) {$N = c^2 + d^2$};

  % Arrow to GCD
  \draw[->,thick] (rep1.south) -- (0,0.5);
  \draw[->,thick] (rep2.south) -- (0,0.5);

  \node[draw, fill=goldaccent!20, rounded corners, minimum width=4cm,
    font=\bfseries] (gcd) at (0,0) {$\gcd(ac + bd,\; N)$};

  % Factor
  \draw[->,very thick, darkgreen] (gcd.south) -- (0,-1.3);
  \node[draw, fill=darkgreen!15, rounded corners, minimum width=2.5cm,
    font=\bfseries\color{darkgreen}] at (0,-1.8) {Factor of $N$!};
\end{tikzpicture}
\end{center}

\section{Road II: Gaussian Composition}

The Pythagorean equation is secretly about norms in the Gaussian integers
$\Z[i]$. If $a^2 + b^2 = N$, then $N = (a+bi)(a-bi) = |a+bi|^2$ in $\Z[i]$.

Gaussian composition---multiplying elements of $\Z[i]$---preserves the
``Pythagorean property,'' and the factorization structure of $N$ in $\Z[i]$
reveals the factorization of $N$ in $\Z$.

\section{Road III: The Tree Sieve}

The Berggren tree provides a systematic \emph{sieve}: by traversing the
tree and collecting triples whose hypotenuse is ``smooth'' (has only small
prime factors), we accumulate enough relations to factor $N$ by
linear algebra modulo 2---the same strategy used by the quadratic sieve
and the number field sieve.

\begin{insightbox}[Three Roads, One Destination]
All three methods---Euler's dual representation, Gaussian composition,
and the tree sieve---ultimately rely on the same algebraic identity:
$(c-b)(c+b) = a^2$. The difference is in how they \emph{search} for
suitable triples.
\end{insightbox}

% ───────────────────────────────────────────────────────────────────
\chapter{Publication-Quality Proofs}
\label{ch:paper-proofs}

\epigraph{\itshape ``Beware of bugs in the above code;\\
I have only proved it correct, not tested it.''}
{--- Donald Knuth}

\lettrine[lines=3]{T}{his chapter} presents the core theorems of the
Berggren--Lorentz correspondence in publication-quality form, suitable
for a refereed mathematics journal. Every theorem has been formally
verified with a clean axiom audit: no \sorry{}, no custom axioms,
only the standard axioms of Lean's type theory (\texttt{propext},
\texttt{Classical.choice}, \texttt{Quot.sound}).

\section{The Eight Main Theorems}

\begin{theorem}[Lorentz Preservation, Thm.~3.1]
For each $X \in \{A, B, C\}$: $B_X^\top\, Q\, B_X = Q$ where
$Q = \mathrm{diag}(1,1,-1)$.
\end{theorem}

\begin{theorem}[Pythagorean Preservation, Thm.~3.2]
If $(a,b,c)$ is Pythagorean and $M$ is a Berggren matrix, then
$M \cdot (a,b,c)^\top$ is Pythagorean.
\end{theorem}

\begin{theorem}[Tree Soundness, Thm.~3.3]
Every triple reachable from $(3,4,5)$ via finite applications of
$B_A, B_B, B_C$ is a Pythagorean triple.
\end{theorem}

\begin{theorem}[Factoring Identity, Thm.~3.4]
For any Pythagorean triple: $(c-b)(c+b) = a^2$.
\end{theorem}

\begin{theorem}[Euclid Parametrization, Thm.~3.5]
$(m^2-n^2)^2 + (2mn)^2 = (m^2+n^2)^2$ for all integers $m, n$.
\end{theorem}

\begin{theorem}[Pell Recurrence, Thm.~4.1]
B-branch hypotenuses satisfy $c_{n+2} = 6c_{n+1} - c_n$.
\end{theorem}

\begin{theorem}[Determinants]
$\det(B_A) = 1$, $\det(B_B) = -1$, $\det(B_C) = 1$.
\end{theorem}

\begin{theorem}[A-branch Descent, Thm.~4.4]
Consecutive Euclid parameters $(m, m-1)$ descend by pure $A$-steps.
\end{theorem}

\begin{center}
\begin{tikzpicture}[scale=0.8]
  % Theorem dependency graph
  \node[draw, rounded corners, fill=deepblue!15, font=\small\bfseries,
    minimum width=2.8cm, text width=2.2cm, align=center] (t31) at (0,5) {Thm 3.1\\Lorentz};
  \node[draw, rounded corners, fill=deepblue!15, font=\small\bfseries,
    minimum width=2.8cm, text width=2.2cm, align=center] (t32) at (4,5) {Thm 3.2\\Pythag.};
  \node[draw, rounded corners, fill=lorentzred!15, font=\small\bfseries,
    minimum width=2.8cm, text width=2.2cm, align=center] (t33) at (2,3) {Thm 3.3\\Soundness};
  \node[draw, rounded corners, fill=darkgreen!15, font=\small\bfseries,
    minimum width=2.8cm, text width=2.2cm, align=center] (t34) at (-2,3) {Thm 3.4\\Factoring};
  \node[draw, rounded corners, fill=goldaccent!15, font=\small\bfseries,
    minimum width=2.8cm, text width=2.2cm, align=center] (t35) at (6,3) {Thm 3.5\\Euclid};
  \node[draw, rounded corners, fill=fermatpurple!15, font=\small\bfseries,
    minimum width=2.8cm, text width=2.2cm, align=center] (t41) at (0,1) {Thm 4.1\\Pell};
  \node[draw, rounded corners, fill=cayleyorange!15, font=\small\bfseries,
    minimum width=2.8cm, text width=2.2cm, align=center] (t44) at (4,1) {Thm 4.4\\Descent};

  \draw[->,thick] (t31) -- (t33);
  \draw[->,thick] (t32) -- (t33);
  \draw[->,thick] (t33) -- (t34);
  \draw[->,thick] (t35) -- (t33);
  \draw[->,thick] (t33) -- (t41);
  \draw[->,thick] (t35) -- (t44);

  \node[font=\small\itshape] at (2,-0.5) {Theorem dependency graph};
\end{tikzpicture}
\end{center}

% ═══════════════════════════════════════════════════════════════════
%  PART II: THE CHANNELS
% ═══════════════════════════════════════════════════════════════════
\part{The Channels of Number}

% ───────────────────────────────────────────────────────────────────
\chapter{Higher $k$-Tuple Factoring}
\label{ch:k-tuple}

\epigraph{\itshape ``The art of doing mathematics consists in finding\\
that special case which contains all the germs of generality.''}
{--- David Hilbert}

\lettrine[lines=3]{W}{hy stop} at triples? If three squares can sum
to a square---$a^2 + b^2 = c^2$---then surely four squares can sum to a
square, and five, and six. Each additional ``channel'' provides new
information about the number being factored.

\section{The Channel Framework}

A \emph{Pythagorean $k$-tuple} is a tuple $(x_1, x_2, \ldots, x_k)$ with
\[
x_1^2 + x_2^2 + \cdots + x_{k-1}^2 = x_k^2.
\]

For $k = 4$ (quadruples), each subset of three terms gives a ``channel'':
\[
\text{Channel}_i = x_k^2 - x_i^2 = \sum_{j \neq i, k} x_j^2.
\]

A quadruple $(a,b,c,d)$ with $a^2 + b^2 + c^2 = d^2$ gives three
primary channels:
\begin{align*}
\text{Channel}_a &= b^2 + c^2 = d^2 - a^2 = (d-a)(d+a),\\
\text{Channel}_b &= a^2 + c^2 = d^2 - b^2 = (d-b)(d+b),\\
\text{Channel}_c &= a^2 + b^2 = d^2 - c^2 = (d-c)(d+c).
\end{align*}

Each channel provides a \emph{difference of squares} factorization.
For octuplets ($k = 8$), there are 7 primary channels and
$\binom{7}{2} = 21$ pairwise channels---28 independent attempts
to factor.

\begin{center}
\begin{tikzpicture}[scale=0.85]
  % Central node
  \node[draw, circle, fill=goldaccent!20, minimum size=1.2cm,
    font=\small\bfseries] (N) at (0,0) {$N$};

  % Channel nodes (k=4: 3 channels)
  \foreach \i/\lbl/\col/\ang in {
    1/{Ch$_a$}/deepblue/90,
    2/{Ch$_b$}/lorentzred/210,
    3/{Ch$_c$}/darkgreen/330} {
    \node[draw, rounded corners, fill=\col!15, minimum width=1.5cm,
      font=\small] (c\i) at (\ang:2.5) {\lbl};
    \draw[->,thick,\col] (N) -- (c\i);
  }

  % k=8: additional channels
  \foreach \i/\ang in {4/30,5/150,6/270} {
    \node[draw, rounded corners, fill=gray!15, minimum width=1cm,
      font=\tiny] (c\i) at (\ang:3.8) {Ch$_\i$};
    \draw[->,thin,gray] (N) -- (c\i);
  }

  \node[font=\small\itshape] at (0,-4.5) {Multi-channel factoring: each channel gives $N$ as $(d-x)(d+x)$};

  % Label
  \node[font=\footnotesize, text width=3cm, align=center] at (5,0)
    {$k=4$: 3 channels\\$k=8$: 28 channels\\More channels $=$\\more chances};
\end{tikzpicture}
\end{center}

\section{Cross-Dimensional Lifting}

A Pythagorean triple $(a,b,c)$ can be \emph{lifted} to a quadruple by
finding $d$ such that $a^2 + b^2 + c^2 = d^2$ (i.e., $c^2 + c^2 = d^2 - a^2 + b^2$, etc.).
The quaternion norm multiplicativity---Euler's four-square identity---guarantees
that products of representable numbers remain representable:
\[
(x_1^2 + x_2^2 + x_3^2 + x_4^2)(y_1^2 + y_2^2 + y_3^2 + y_4^2)
= z_1^2 + z_2^2 + z_3^2 + z_4^2.
\]

\section{The Reflection Operator}

For quadruples, the matrix
\[
R_{1111} = \frac{1}{2}\begin{pmatrix} 1&1&1&1\\1&1&-1&-1\\1&-1&1&-1\\1&-1&-1&1\end{pmatrix}
\]
preserves the null cone $a^2+b^2+c^2 = d^2$ (when the triple has the
right parity) and reveals factor structure through the transformed
coordinates.

% ───────────────────────────────────────────────────────────────────
\chapter{Quantum Grover Acceleration}
\label{ch:quantum}

\epigraph{\itshape ``Nature isn't classical, dammit, and if you want\\
to make a simulation of nature, you'd better make it\\
quantum mechanical.''}
{--- Richard Feynman}

\lettrine[lines=3]{C}{an quantum} computers break the $\sqrt{N}$ barrier?
Grover's algorithm---the quantum analogue of exhaustive search---provides
a quadratic speedup for unstructured search problems. Applied to
Berggren tree descent, it reduces the factoring complexity from
$O(\sqrt{N})$ to $O(N^{1/4})$.

\section{The Oracle Construction}

Grover's algorithm requires an \emph{oracle} that recognizes solutions.
For tree factoring, the oracle checks: ``Does inverse Berggren descent
from this triple reach the root?'' The key insight is that tree descent
is \emph{deterministic}---at most one inverse branch produces a valid
triple at each step.

\begin{theorem}[Deterministic Descent]
For any Pythagorean triple $(a,b,c)$ with $c > 0$, at most one of the
three inverse Berggren matrices $B_A^{-1}$, $B_B^{-1}$, $B_C^{-1}$
produces a triple with all positive coordinates.
\end{theorem}

This determinism is crucial: it means each oracle call takes $O(\sqrt{N})$
classical steps, and Grover's algorithm makes $O(N^{1/4})$ oracle calls,
for a total of $O(N^{1/4} \cdot \sqrt{N}) = O(N^{3/4})$---wait, that's
worse! The trick is to embed the descent \emph{inside} the Grover oracle,
checking only the first few steps.

\begin{center}
\begin{tikzpicture}[scale=0.9]
  % Quantum circuit
  \foreach \y in {0,0.8,1.6,2.4} {
    \draw[thick] (0,\y) -- (8,\y);
  }

  % Qubit labels
  \node[font=\small] at (-0.5,2.4) {$|0\rangle$};
  \node[font=\small] at (-0.5,1.6) {$|0\rangle$};
  \node[font=\small] at (-0.5,0.8) {$|0\rangle$};
  \node[font=\small] at (-0.5,0) {$|0\rangle$};

  % Hadamard gates
  \foreach \y in {0,0.8,1.6,2.4} {
    \node[draw, fill=deepblue!15, minimum size=0.5cm, font=\small] at (1,\y) {$H$};
  }

  % Oracle
  \node[draw, fill=goldaccent!20, minimum width=1.2cm, minimum height=3.2cm,
    font=\small\bfseries] at (3,1.2) {$U_f$};

  % Diffusion
  \node[draw, fill=lorentzred!15, minimum width=1.2cm, minimum height=3.2cm,
    font=\small\bfseries] at (5,1.2) {$D$};

  % Measurement
  \foreach \y in {0,0.8,1.6,2.4} {
    \node[draw, fill=darkgreen!15, minimum size=0.5cm, font=\small] at (7,\y) {$M$};
  }

  % Repeat bracket
  \draw[decorate, decoration={brace, amplitude=8pt}, thick]
    (2.2,3.2) -- (5.8,3.2) node[midway, above=10pt, font=\small] {repeat $O(N^{1/4})$ times};

  \node[font=\small\itshape] at (4,-1) {Grover's algorithm applied to Berggren tree search};
\end{tikzpicture}
\end{center}

% ───────────────────────────────────────────────────────────────────
\chapter{Complexity Bounds, Proven}
\label{ch:complexity}

\epigraph{\itshape ``In theory, there is no difference between theory and practice.\\
In practice, there is.''}
{--- Yogi Berra}

\lettrine[lines=3]{T}{he central} complexity result of this work is
machine-verified: \emph{Pythagorean tree factoring for balanced semiprimes
$N = pq$ requires $\Theta(\sqrt{N})$ operations.}

\section{The Upper Bound}

For $N = pq$ with $p \leq q$, the tree descent requires at most
$O(m)$ steps where $m \sim \sqrt{N}$ is the larger Euclid parameter.
Since each step is a constant number of arithmetic operations on
$O(\log N)$-bit numbers, the total work is $O(\sqrt{N} \cdot \mathrm{polylog}(N))$.

\section{The Lower Bound}

The matching lower bound follows from the Lattice--Tree Correspondence
(Chapter~\ref{ch:lattice-tree}): since descent is equivalent to the
Euclidean algorithm on $(m,n)$, and consecutive parameters $(m, m-1)$
force $m - 2$ steps of pure $A$-descent, we get $\Omega(\sqrt{N})$
in the worst case.

\begin{center}
\begin{tikzpicture}[scale=0.85]
  \begin{axis}[
    xlabel={$\log_2 N$},
    ylabel={Steps},
    xmin=0, xmax=60,
    ymin=0, ymax=1.2e9,
    legend pos=north west,
    grid=major,
    width=10cm, height=6cm,
    scaled y ticks=false,
    yticklabel style={/pgf/number format/fixed},
  ]
    % sqrt(N) curve
    \addplot[deepblue, very thick, domain=1:60, samples=50]
      {2^(x/2)};
    \addlegendentry{$\sqrt{N}$ (tree descent)}

    % N^{1/4} curve
    \addplot[lorentzred, very thick, dashed, domain=1:60, samples=50]
      {2^(x/4)};
    \addlegendentry{$N^{1/4}$ (Grover)}

    % Trial division
    \addplot[gray, thick, dotted, domain=1:60, samples=50]
      {2^(x/2)};
    \addlegendentry{$\sqrt{N}$ (trial division)}

  \end{axis}
\end{tikzpicture}
\end{center}

% ═══════════════════════════════════════════════════════════════════
%  PART III: THE ALGEBRA
% ═══════════════════════════════════════════════════════════════════
\part{The Algebraic Cosmos}

% ───────────────────────────────────────────────────────────────────
\chapter{The Cayley--Dickson Hierarchy}
\label{ch:cayley-dickson}

\epigraph{\itshape ``The imaginary number is a fine and wonderful recourse\\
of the divine spirit, almost an amphibian between being and not being.''}
{--- Gottfried Wilhelm Leibniz}

\lettrine[lines=3]{T}{here is} a ladder of number systems, and at each
rung you gain a new power but lose an old one. The Cayley--Dickson
construction builds each system by \emph{doubling} the previous one,
and at each doubling, an algebraic property is sacrificed on the altar
of higher dimension.

\section{The Four Channels}

\begin{center}
\begin{tikzpicture}[
  box/.style={draw, rounded corners, minimum width=3cm, minimum height=1.2cm,
    font=\small, text width=2.8cm, align=center},
  arrow/.style={->, thick, font=\small},
  loss/.style={font=\small\itshape\color{lorentzred}, text width=3cm, align=center},
  gain/.style={font=\small\itshape\color{darkgreen}, text width=3cm, align=center}
]
  \node[box, fill=deepblue!10] (R) at (0,6) {$\R$\\Reals\\dim = 1};
  \node[box, fill=deepblue!15] (C) at (0,4) {$\C$\\Complex\\dim = 2};
  \node[box, fill=deepblue!20] (H) at (0,2) {$\HH$\\Quaternions\\dim = 4};
  \node[box, fill=deepblue!25] (O) at (0,0) {$\OO$\\Octonions\\dim = 8};
  \node[box, fill=lorentzred!15] (S) at (0,-2) {$\mathbb{S}$\\Sedenions\\dim = 16};

  \draw[arrow] (R) -- (C);
  \draw[arrow] (C) -- (H);
  \draw[arrow] (H) -- (O);
  \draw[arrow, lorentzred] (O) -- (S);

  % Losses
  \node[loss] at (3.8,5) {Lose:\\total ordering};
  \node[loss] at (3.8,3) {Lose:\\commutativity\\$ij \neq ji$};
  \node[loss] at (3.8,1) {Lose:\\associativity\\$(ij)k \neq i(jk)$};
  \node[loss] at (3.8,-1) {Lose:\\division!\\zero divisors appear};

  % Gains
  \node[gain] at (-3.8,5) {Gain:\\algebraic closure};
  \node[gain] at (-3.8,3) {Gain:\\3D rotations};
  \node[gain] at (-3.8,1) {Gain:\\$E_8$ lattice};
  \node[gain] at (-3.8,-1) {Gain: \ldots \\nothing useful};

  % Channel boundary marker
  \draw[lorentzred, very thick, dashed] (-5.5,-1) -- (5.5,-1);
  \node[lorentzred, font=\bfseries] at (5.5,-1.5) {CHANNEL BREAKS};
\end{tikzpicture}
\end{center}

\section{Composition Algebras and Factoring}

At each level, there is a corresponding \emph{composition identity}:
\begin{itemize}
\item \textbf{$\R \to \C$:} Brahmagupta--Fibonacci identity (2 squares):
$(a^2+b^2)(c^2+d^2) = (ac-bd)^2 + (ad+bc)^2$
\item \textbf{$\C \to \HH$:} Euler's four-square identity (4 squares)
\item \textbf{$\HH \to \OO$:} Degen's eight-square identity (8 squares)
\end{itemize}

Each identity says that the product of two sums of $k$ squares is itself
a sum of $k$ squares, with explicit formulas for the resulting terms.
These identities are the algebraic engines that power multi-channel factoring.

\begin{theorem}[Machine-Verified]
Quaternion multiplication is \emph{not} commutative: there exist
$a, b \in \HH$ with $ab \neq ba$.

In Lean 4: $a = (0,1,0,0)$, $b = (0,0,1,0)$. Then $ab$ has
$\mathrm{im}_K = 1$ while $ba$ has $\mathrm{im}_K = -1$.
\end{theorem}

% ───────────────────────────────────────────────────────────────────
\chapter{Fermat's Last Theorem: What Fits in the Margin?}
\label{ch:fermat}

\epigraph{\itshape ``I have discovered a truly marvelous demonstration of this\\
proposition that this margin is too narrow to contain.''}
{--- Pierre de Fermat, 1637}

\lettrine[lines=3]{F}{ermat's} Last Theorem is the most famous ``theorem''
in mathematics---famous in part because for 358 years it was not actually a
theorem but a \emph{conjecture}. When Andrew Wiles finally proved it in 1995,
he needed 109 pages of the most sophisticated mathematics of the 20th century.

But the cases $n = 3$ and $n = 4$ really do fit in a margin. And they
are formally verified in Lean.

\section{The Case $n = 4$: Infinite Descent}

Fermat himself proved this case, using his method of infinite descent:
assume a solution exists with minimal $c$, then construct a strictly
smaller solution, contradicting minimality.

\begin{theorem}[Fermat, $\sim$1640; Machine-Verified]
There are no positive integers $a, b, c$ with $a^4 + b^4 = c^4$.

In fact, the stronger statement holds: $a^4 + b^4 \neq c^2$ for
positive $a, b, c$.
\end{theorem}

\section{The Case $n = 3$: Euler's Proof}

Euler proved this in 1770, using factorization in $\Z[\omega]$ where
$\omega = e^{2\pi i/3}$ is a primitive cube root of unity.

\begin{theorem}[Euler, 1770; Machine-Verified]
There are no positive integers $a, b, c$ with $a^3 + b^3 = c^3$.
\end{theorem}

\section{The Full Theorem}

The full Fermat's Last Theorem requires the modularity theorem
(Taniyama--Shimura--Weil) and is not yet fully formalized in Lean,
though a major formalization effort is underway.

\begin{center}
\begin{tikzpicture}[scale=0.85]
  % Timeline
  \draw[->,very thick] (0,0) -- (12,0);

  % Events
  \foreach \x/\yr/\evt in {
    0.5/1637/Fermat's\\note,
    2.5/1770/Euler\\$n{=}3$,
    4/1825/Dirichlet\\$n{=}5$,
    5.5/1847/Kummer\\regular,
    8/1995/Wiles\\proof,
    10.5/2024/Lean\\$n{=}3{,}4$} {
    \fill (\x,0) circle(3pt);
    \node[font=\tiny, text width=1.2cm, align=center, above] at (\x,0.3) {\evt};
    \node[font=\tiny, below] at (\x,-0.2) {\yr};
  }

  \node[font=\small\itshape] at (6,-1.2) {358 years from conjecture to proof};
\end{tikzpicture}
\end{center}

\section{Why the Margin Was Too Small}

The most likely explanation for Fermat's ``marvelous proof'' is that he
attempted factorization in cyclotomic integers $\Z[\zeta_n]$ and assumed
unique factorization---which fails for $n = 23$ and many other values.
Kummer discovered this obstruction in 1847, two centuries after Fermat.

The margin was not too small. The proof was too big. And it took the
collective genius of Frey, Serre, Ribet, and Wiles to find it.

% ───────────────────────────────────────────────────────────────────
\chapter{Congruence of Squares}
\label{ch:congruence}

\epigraph{\itshape ``God does arithmetic.''}
{--- Carl Friedrich Gauss}

\lettrine[lines=3]{E}{very modern} sub-exponential factoring algorithm---the
quadratic sieve, the number field sieve, the elliptic curve method---rests
on a single observation:

\begin{insightbox}[The Congruence of Squares Principle]
If $x^2 \equiv y^2 \pmod{N}$ with $x \not\equiv \pm y \pmod{N}$,
then $\gcd(x - y, N)$ is a nontrivial factor of $N$.
\end{insightbox}

This is because $N \mid (x^2 - y^2) = (x-y)(x+y)$, but $N \nmid (x-y)$
and $N \nmid (x+y)$. So $N$ must split across the two factors, and
$\gcd(x-y, N)$ captures one of them.

\begin{center}
\begin{tikzpicture}[scale=0.9]
  % The factoring chain
  \node[draw, rounded corners, fill=deepblue!10, text width=3cm, align=center]
    (start) at (0,3) {Find $x, y$\\$x^2 \equiv y^2$\\$\pmod{N}$};
  \node[draw, rounded corners, fill=goldaccent!15, text width=3cm, align=center]
    (check) at (0,1) {Check\\$x \not\equiv \pm y$\\$\pmod{N}$};
  \node[draw, rounded corners, fill=darkgreen!20, text width=3cm, align=center,
    font=\bfseries] (win) at (0,-1) {$\gcd(x-y, N)$\\= factor!};

  \draw[->,thick] (start) -- (check);
  \draw[->,thick] (check) -- (win);

  % Methods
  \node[draw, rounded corners, fill=lorentzred!10, font=\small] (qs) at (5,3) {Quadratic Sieve};
  \node[draw, rounded corners, fill=lorentzred!10, font=\small] (nfs) at (5,1.5) {Number Field Sieve};
  \node[draw, rounded corners, fill=lorentzred!10, font=\small] (ecm) at (5,0) {Elliptic Curve Method};
  \node[draw, rounded corners, fill=pythagorean!10, font=\small] (tree) at (5,-1.5) {Berggren Tree};

  \draw[->,thick,gray] (qs.west) -- (start.east);
  \draw[->,thick,gray] (nfs.west) -- (start.east);
  \draw[->,thick,gray] (ecm.west) -- (check.east);
  \draw[->,thick,gray] (tree.west) -- (check.east);

  \node[font=\small\itshape, text width=3.5cm, align=center] at (5,-3)
    {All these methods search\\for congruent squares};
\end{tikzpicture}
\end{center}

The connection to Pythagorean triples is immediate: $a^2 + b^2 = c^2$
gives $(c-b)(c+b) = a^2$, so $c^2 \equiv b^2 \pmod{a^2}$.

% ═══════════════════════════════════════════════════════════════════
%  PART IV: THE FRONTIERS
% ═══════════════════════════════════════════════════════════════════
\part{The Frontiers}

% ───────────────────────────────────────────────────────────────────
\chapter{Quadruple Factor Theory}
\label{ch:quadruple}

\epigraph{\itshape ``Mathematics is the queen of the sciences and\\
number theory is the queen of mathematics.''}
{--- Carl Friedrich Gauss}

\lettrine[lines=3]{P}{ythagorean quadruples}---solutions to
$a^2 + b^2 + c^2 = d^2$---provide a richer factoring landscape than triples.
The ``Shared Factor Bridge'' theorem shows that when $d$ is composite,
the channel structure encodes information about its factors.

\section{The Shared Factor Bridge}

\begin{theorem}[Shared Factor Bridge]
If $g$ divides both channels $(a^2+b^2)$ and $(a^2+c^2)$, then $g$
divides $b^2 - c^2 = (b-c)(b+c)$.
\end{theorem}

This means that pairwise GCDs of channel values reveal factor structure.
For a composite hypotenuse $d = pq$, the channels are constrained by
the factorization, and cross-channel GCDs extract factors.

\section{Complete Parametrization}

Every primitive Pythagorean quadruple can be written as:
\[
a = m^2 + n^2 - p^2 - q^2,\quad b = 2(mq+np),\quad
c = 2(nq - mp),\quad d = m^2 + n^2 + p^2 + q^2
\]
for integers $m, n, p, q$ with $\gcd(m,n,p,q) = 1$ and appropriate
parity conditions.

\begin{center}
\begin{tikzpicture}[scale=0.8]
  % 3D-ish representation of quadruple
  \coordinate (O) at (0,0);
  \coordinate (A) at (3,0);
  \coordinate (B) at (1.5,2);
  \coordinate (C) at (0.5,1);
  \coordinate (D) at (2.5,3);

  % Faces
  \fill[deepblue!10] (O) -- (A) -- (D) -- (B) -- cycle;
  \fill[lorentzred!10] (O) -- (B) -- (D) -- (C) -- cycle;
  \fill[darkgreen!10] (O) -- (A) -- (D) -- (C) -- cycle;

  \draw[thick] (O) -- (A) node[midway, below, font=\small] {$a$};
  \draw[thick] (O) -- (B) node[midway, left, font=\small] {$b$};
  \draw[thick] (O) -- (C) node[midway, below left, font=\small] {$c$};
  \draw[thick, dashed, lorentzred] (O) -- (D) node[midway, above right, font=\small\color{lorentzred}] {$d$};
  \draw[thin] (A) -- (D);
  \draw[thin] (B) -- (D);
  \draw[thin] (C) -- (D);

  \node[font=\small\itshape] at (1.5,-1) {$a^2 + b^2 + c^2 = d^2$};

  % Channel labels
  \node[deepblue, font=\tiny] at (4.5,1.5) {Ch$_a = d^2-a^2$};
  \node[lorentzred, font=\tiny] at (4.5,0.8) {Ch$_b = d^2-b^2$};
  \node[darkgreen, font=\tiny] at (4.5,0.1) {Ch$_c = d^2-c^2$};
\end{tikzpicture}
\end{center}

% ───────────────────────────────────────────────────────────────────
\chapter{GCD Cascades and Factor Extraction}
\label{ch:gcd-cascade}

\epigraph{\itshape ``The greatest common divisor is the most democratic\\
of all mathematical operations.''}
{--- Paul Klemstine}

\lettrine[lines=3]{W}{hen a number} has \emph{multiple} representations
as a sum of squares, the interplay between these representations creates
a cascade of GCD computations, each potentially revealing a new factor.

\section{The Channel GCD Lattice}

Given a Pythagorean quadruple $(a,b,c,d)$, the three channels generate
pairwise GCDs:
\[
g_{ab} = \gcd(d^2-a^2,\; d^2-b^2),\quad
g_{ac} = \gcd(d^2-a^2,\; d^2-c^2),\quad
g_{bc} = \gcd(d^2-b^2,\; d^2-c^2).
\]

These GCDs satisfy a lattice structure: if $g$ divides all three channels,
it divides all pairwise squared differences $a^2-b^2$, $a^2-c^2$, $b^2-c^2$.

\section{Multi-Representation Cascades}

When a number $N$ has $k$ different representations as a sum of three
squares, we get $\binom{k}{2}$ cross-representation GCDs. Each one is
an independent attempt to factor $N$. As $k$ grows, the probability of
finding a factor approaches 1 exponentially.

\begin{center}
\begin{tikzpicture}[scale=0.8]
  % GCD cascade diagram
  \node[draw, fill=deepblue!15, rounded corners, minimum width=2.5cm] (r1) at (-2,3) {Rep 1};
  \node[draw, fill=deepblue!15, rounded corners, minimum width=2.5cm] (r2) at (2,3) {Rep 2};
  \node[draw, fill=deepblue!15, rounded corners, minimum width=2.5cm] (r3) at (6,3) {Rep 3};

  \node[draw, fill=goldaccent!20, rounded corners] (g12) at (0,1) {$\gcd$};
  \node[draw, fill=goldaccent!20, rounded corners] (g13) at (2,0) {$\gcd$};
  \node[draw, fill=goldaccent!20, rounded corners] (g23) at (4,1) {$\gcd$};

  \draw[->,thick] (r1) -- (g12);
  \draw[->,thick] (r2) -- (g12);
  \draw[->,thick] (r1) -- (g13);
  \draw[->,thick] (r3) -- (g13);
  \draw[->,thick] (r2) -- (g23);
  \draw[->,thick] (r3) -- (g23);

  \node[draw, fill=darkgreen!20, rounded corners, font=\bfseries] (factor) at (2,-1.5) {Factor!};
  \draw[->,thick,darkgreen] (g12) -- (factor);
  \draw[->,thick,darkgreen] (g13) -- (factor);
  \draw[->,thick,darkgreen] (g23) -- (factor);

  \node[font=\small\itshape] at (2,-3) {GCD cascade: $\binom{k}{2}$ independent factoring attempts};
\end{tikzpicture}
\end{center}

% ───────────────────────────────────────────────────────────────────
\chapter{Pythagorean Tree Factoring Core}
\label{ch:tree-core}

\epigraph{\itshape ``The best way to predict the future is to implement it.''}
{--- Alan Kay}

\lettrine[lines=3]{T}{his chapter} presents the algorithmic core of
Pythagorean tree factoring: the concrete data structures, the traversal
strategies, and the factor extraction mechanisms.

\section{The Factoring Algorithm}

Given $N$ to factor:
\begin{enumerate}
\item \textbf{Representation:} Find a Pythagorean triple $(a,b,c)$
with $c$ related to $N$ (e.g., $N = c$ or $N \mid c^2$).
\item \textbf{Descent:} Apply inverse Berggren matrices to descend the tree.
\item \textbf{Extraction:} At each step, compute $\gcd(c-b, N)$ and
$\gcd(c+b, N)$.
\item \textbf{Success:} If a nontrivial GCD is found, we have a factor.
\end{enumerate}

\section{Descent Termination}

The key to the algorithm's correctness is that the hypotenuse
\emph{strictly decreases} at each descent step. Since hypotenuses are
positive integers, the descent must terminate---either at the root
$(3,4,5)$ or at a step where a factor is found.

\begin{center}
\begin{tikzpicture}[scale=0.8]
  % Descent visualization
  \foreach \i/\c/\y in {1/1073/5, 2/533/4, 3/265/3, 4/130/2, 5/65/1, 6/5/0} {
    \node[draw, rounded corners, fill=deepblue!\i 0, minimum width=2.5cm,
      font=\small] at (0,\y) {$c = \c$};
  }
  \foreach \y in {4.5,3.5,2.5,1.5,0.5} {
    \draw[->,thick] (0,\y+0.2) -- (0,\y-0.15);
  }

  \node[font=\small\itshape] at (0,-1) {Hypotenuse strictly decreases};
  \node[darkgreen, font=\bfseries] at (3,0) {Root!};
  \draw[->,thick,darkgreen] (1.3,0) -- (2.4,0);
\end{tikzpicture}
\end{center}

% ───────────────────────────────────────────────────────────────────
\chapter{Tropical Geometry Foundations}
\label{ch:tropical}

\epigraph{\itshape ``In the tropical world, addition becomes minimum\\
and multiplication becomes addition.''}
{--- Bernd Sturmfels}

\lettrine[lines=3]{T}{ropical geometry} is what happens when you replace
ordinary arithmetic with the ``min-plus'' algebra: addition becomes
$\min(a,b)$ and multiplication becomes $a + b$. This seemingly absurd
redefinition turns out to be extraordinarily powerful, connecting
algebraic geometry to optimization, phylogenetics, and shortest-path
algorithms.

\section{The Tropical Semiring}

The tropical semiring is $(\R \cup \{+\infty\},\; \oplus,\; \odot)$ where:
\begin{align*}
a \oplus b &= \min(a, b) \quad \text{(tropical addition)},\\
a \odot b &= a + b \quad \text{(tropical multiplication)}.
\end{align*}

\begin{center}
\begin{tikzpicture}[scale=0.85]
  % Classical vs Tropical
  \node[draw, rounded corners, fill=deepblue!10, minimum width=4cm,
    minimum height=2.5cm] at (-3,0) {};
  \node[font=\bfseries, deepblue] at (-3,1.5) {Classical};
  \node[font=\small] at (-3,0.6) {$a + b = $ sum};
  \node[font=\small] at (-3,-0.1) {$a \times b = $ product};
  \node[font=\small] at (-3,-0.8) {$0$ is additive identity};

  \node[draw, rounded corners, fill=tropicalgreen!10, minimum width=4cm,
    minimum height=2.5cm] at (3,0) {};
  \node[font=\bfseries, tropicalgreen] at (3,1.5) {Tropical};
  \node[font=\small] at (3,0.6) {$a \oplus b = \min(a,b)$};
  \node[font=\small] at (3,-0.1) {$a \odot b = a + b$};
  \node[font=\small] at (3,-0.8) {$\infty$ is additive identity};

  \draw[->,very thick, goldaccent] (-0.5,0) -- (0.5,0)
    node[midway, above, font=\small\bfseries\color{goldaccent}] {tropicalize};

  % Tropical line (piecewise linear)
  \begin{scope}[shift={(0,-4)}, scale=0.7]
    \draw[->,thick] (-3,0) -- (4,0) node[right,font=\small] {$x$};
    \draw[->,thick] (0,-1) -- (0,4) node[above,font=\small] {$y$};

    % Tropical line: min(x, y, 0) = 0
    \draw[tropicalgreen, very thick] (-2.5,0) -- (0,0) -- (2.5,2.5);
    \draw[tropicalgreen, very thick] (0,0) -- (0,3);

    \node[font=\small\itshape] at (0,-2) {A tropical ``line'': piecewise linear!};
  \end{scope}
\end{tikzpicture}
\end{center}

The axioms are machine-verified:
\begin{itemize}
\item Commutativity: $\min(a,b) = \min(b,a)$
\item Associativity: $\min(\min(a,b),c) = \min(a,\min(b,c))$
\item Identity: $\min(a, +\infty) = a$
\item Distributivity: $a + \min(b,c) = \min(a+b, a+c)$
\end{itemize}

\section{The Bellman Equation}

The shortest-path problem is naturally tropical: the Bellman equation
$d(v) = \min_{(u,v) \in E} \{d(u) + w(u,v)\}$
is just saying $d(v) = \bigoplus_{u} d(u) \odot w(u,v)$---a tropical
matrix-vector product.

\section{Newton Polygons}

A tropical polynomial $p(x) = \bigoplus_i a_i \odot x^{\odot i} = \min_i(a_i + ix)$
is a piecewise linear function. Its ``roots'' are the slopes where the
minimum switches from one term to another---exactly the slopes of the
Newton polygon of the corresponding classical polynomial. This connects
tropical geometry to $p$-adic valuation theory.

% ───────────────────────────────────────────────────────────────────
\chapter{The Lorentz Group Structure}
\label{ch:lorentz}

\epigraph{\itshape ``Henceforth space by itself, and time by itself, are\\
doomed to fade away into mere shadows, and only a kind\\
of union of the two will preserve an independent reality.''}
{--- Hermann Minkowski}

\lettrine[lines=3]{W}{e return} at last to the Lorentz group, now with the
full force of our algebraic machinery. The Berggren tree is not merely
\emph{related to} the Lorentz group---it \emph{is} a discrete subgroup,
tiling the hyperbolic plane as systematically as wallpaper tiles a wall.

\section{The Tiling Theorem}

Each Berggren matrix preserves the quadratic form
$\lorentzQ(a,b,c) = a^2 + b^2 - c^2$, as we proved in
Chapter~\ref{ch:berggren-lorentz}. The three matrices $B_A, B_B, B_C$
generate a free subgroup of $O(2,1;\Z)$ that acts freely and properly
discontinuously on the upper sheet of the hyperboloid $\lorentzQ = -1$.

\begin{center}
\begin{tikzpicture}[scale=1.0]
  % Hyperboloid sheet (upper)
  \draw[domain=-2:2, smooth, variable=\x, lorentzred, very thick]
    plot ({\x}, {sqrt(1 + \x*\x)});
  \draw[domain=-2:2, smooth, variable=\x, lorentzred, very thick, dashed]
    plot ({\x}, {-sqrt(1 + \x*\x)});

  % Light cone
  \draw[gray, thin] (-3,-3) -- (3,3);
  \draw[gray, thin] (-3,3) -- (3,-3);

  % Fundamental domain patches
  \fill[deepblue!20, opacity=0.5] (0,1) -- (0.5,1.12) -- (0.3,1.04) -- cycle;
  \fill[lorentzred!20, opacity=0.5] (0.5,1.12) -- (1.0,1.41) -- (0.7,1.22) -- cycle;
  \fill[darkgreen!20, opacity=0.5] (-0.5,1.12) -- (0,1) -- (-0.3,1.04) -- cycle;

  % Tree points on hyperboloid
  \fill[deepblue] (0,1) circle(2.5pt);
  \fill[deepblue] (0.5,1.12) circle(2pt);
  \fill[deepblue] (-0.5,1.12) circle(2pt);
  \fill[deepblue] (1.0,1.41) circle(1.5pt);
  \fill[deepblue] (1.5,1.80) circle(1.5pt);
  \fill[deepblue] (-1.0,1.41) circle(1.5pt);

  % Labels
  \node[font=\small] at (0,1) [above right] {$(3,4,5)$};
  \node[lorentzred, font=\small] at (3.5,3) {light cone};
  \node[font=\small\itshape] at (0,-3.5) {Berggren tree on the hyperboloid $\lorentzQ = 0$};

  % Axes
  \draw[->,thick] (-3,0) -- (3.5,0) node[right,font=\small] {$a/c$};
  \draw[->,thick] (0,-3) -- (0,3.5) node[above,font=\small] {$b/c$};
\end{tikzpicture}
\end{center}

\section{The Semiprime Counting Theorem}

\begin{theorem}[Semiprime Counting]
For distinct odd primes $p \neq q$, the semiprime $N = pq$ has exactly
4 primitive Pythagorean triples with hypotenuse $N$.

This follows from the divisor formula: the number of representations of
$N$ as $a^2 + b^2$ equals $(d_1(N^2) - 1)/2$ where $d_1$ counts
divisors of $N^2$ less than $N$. For $N = pq$, we have
$d_1((pq)^2) = d_1(p^2 q^2) = 9$, giving $(9-1)/2 = 4$ triples.
\end{theorem}

This is cryptographically significant: an RSA modulus $N = pq$ has
exactly 4 Pythagorean triples, and finding any one of them is
equivalent to factoring $N$.

% ═══════════════════════════════════════════════════════════════════
%  CODA
% ═══════════════════════════════════════════════════════════════════

\chapter*{Coda: The Eternal Braid}
\addcontentsline{toc}{chapter}{Coda: The Eternal Braid}

\lettrine[lines=3]{W}{e began} with a schoolchild's triangle and a
scribble in a margin. We end with a web of connections that spans
millennia of mathematics:

\begin{itemize}
\item Pythagoras's theorem ($\sim$500 BCE) generates an infinite tree.
\item That tree lives in Lorentz's group (1905 CE).
\item Lorentz's group tiles Poincar\'e's disk (1882 CE).
\item The tree traversal is Euclid's algorithm ($\sim$300 BCE).
\item The complexity matches trial division (always known, now proven).
\item The escape route uses Lenstra's LLL algorithm (1982 CE).
\item Grover's quantum algorithm (1996 CE) provides a quadratic speedup.
\item The factoring connection threatens Rivest--Shamir--Adleman (1978 CE).
\item And the Cayley--Dickson hierarchy ($\sim$1900 CE) opens higher channels.
\end{itemize}

Each connection has been machine-verified. The proofs exist not as
arguments to be debated but as \emph{artifacts} to be inspected---digital
objects whose correctness is guaranteed by the same logical foundations
that underpin all of mathematics.

\bigskip

The braid is eternal because it is \emph{necessary}. These connections
are not accidents of notation or artifacts of clever encoding. They exist
because the integers demand them. The Lorentz group preserves the
Pythagorean form because that is what quadratic forms \emph{do}. The tree
matches the Euclidean algorithm because lattice reduction \emph{is} tree
descent. The quantum speedup applies because Grover's algorithm applies
to \emph{all} unstructured search.

And Fermat was right about $a^n + b^n \neq c^n$, even if his proof was
wrong, because mathematics is forgiving of flawed proofs---it is the
theorems that endure.

\bigskip
\begin{center}
\begin{tikzpicture}
  % Three braided strands
  \foreach \i in {0,1,...,20} {
    \pgfmathsetmacro{\x}{\i * 0.5}
    \pgfmathsetmacro{\ya}{0.5*sin(\i * 60)}
    \pgfmathsetmacro{\yb}{0.5*sin(\i * 60 + 120)}
    \pgfmathsetmacro{\yc}{0.5*sin(\i * 60 + 240)}
    \ifnum\i>0
      \pgfmathsetmacro{\px}{(\i-1) * 0.5}
      \pgfmathsetmacro{\pya}{0.5*sin((\i-1) * 60)}
      \pgfmathsetmacro{\pyb}{0.5*sin((\i-1) * 60 + 120)}
      \pgfmathsetmacro{\pyc}{0.5*sin((\i-1) * 60 + 240)}
      \draw[deepblue, thick] (\px,\pya) -- (\x,\ya);
      \draw[lorentzred, thick] (\px,\pyb) -- (\x,\yb);
      \draw[darkgreen, thick] (\px,\pyc) -- (\x,\yc);
    \fi
  }

  \node[deepblue, font=\small] at (11,0.8) {Pythagoras};
  \node[lorentzred, font=\small] at (11,0) {Lorentz};
  \node[darkgreen, font=\small] at (11,-0.8) {Proof};
\end{tikzpicture}
\end{center}

\vfill
\begin{center}
\textit{Soli Deo Gloria}
\end{center}

% ═══════════════════════════════════════════════════════════════════
%  APPENDICES
% ═══════════════════════════════════════════════════════════════════
\appendix
\part*{Appendices}
\addcontentsline{toc}{part}{Appendices}

\chapter{Lean 4 Source Code}
\label{app:lean}

The following pages contain the complete Lean 4 source files for all
machine-verified theorems in this book. Each file is self-contained
(modulo \texttt{import Mathlib}) and can be compiled with Lean 4 and
Mathlib.

\bigskip

\noindent\textbf{Verification Environment:}
\begin{itemize}
\item Lean 4 with Mathlib
\item All theorems checked by Lean's kernel
\item Axiom audit: only \texttt{propext}, \texttt{Classical.choice},
  \texttt{Quot.sound} used
\end{itemize}

\bigskip

% Auto-generated appendix of Lean source files

\section{The Berggren--Lorentz Correspondence}
\label{app:01_BerggrenLorentzCorrespondence}
\texttt{01_BerggrenLorentzCorrespondence.lean}

\lstinputlisting[language=Lean]{../01_BerggrenLorentzCorrespondence.lean}


\section{Lattice--Tree Correspondence}
\label{app:02_LatticeTreeCorrespondence}
\texttt{02_LatticeTreeCorrespondence.lean}

\lstinputlisting[language=Lean]{../02_LatticeTreeCorrespondence.lean}


\section{Hyperbolic Shortcuts for Factoring}
\label{app:03_HyperbolicShortcutsFactoring}
\texttt{03_HyperbolicShortcutsFactoring.lean}

\lstinputlisting[language=Lean]{../03_HyperbolicShortcutsFactoring.lean}


\section{Three Roads from Pythagoras}
\label{app:04_ThreeRoadsFromPythagoras}
\texttt{04_ThreeRoadsFromPythagoras.lean}

\lstinputlisting[language=Lean]{../04_ThreeRoadsFromPythagoras.lean}


\section{Berggren--Lorentz Paper Proofs}
\label{app:05_BerggrenLorentzPaperProofs}
\texttt{05_BerggrenLorentzPaperProofs.lean}

\lstinputlisting[language=Lean]{../05_BerggrenLorentzPaperProofs.lean}


\section{Higher $k$-Tuple Factoring}
\label{app:06_HigherKTupleFactoring}
\texttt{06_HigherKTupleFactoring.lean}

\lstinputlisting[language=Lean]{../06_HigherKTupleFactoring.lean}


\section{Quantum Grover Tree Factoring}
\label{app:07_QuantumGroverTreeFactoring}
\texttt{07_QuantumGroverTreeFactoring.lean}

\lstinputlisting[language=Lean]{../07_QuantumGroverTreeFactoring.lean}


\section{Complexity Bounds, Proven}
\label{app:08_ComplexityBoundsProven}
\texttt{08_ComplexityBoundsProven.lean}

\lstinputlisting[language=Lean]{../08_ComplexityBoundsProven.lean}


\section{Cayley--Dickson Hierarchy}
\label{app:09_CayleyDicksonHierarchy}
\texttt{09_CayleyDicksonHierarchy.lean}

\lstinputlisting[language=Lean]{../09_CayleyDicksonHierarchy.lean}


\section{Fermat's Last Theorem}
\label{app:10_FermatLastTheorem}
\texttt{10_FermatLastTheorem.lean}

\lstinputlisting[language=Lean]{../10_FermatLastTheorem.lean}


\section{Congruence of Squares Factoring}
\label{app:11_CongruenceOfSquaresFactoring}
\texttt{11_CongruenceOfSquaresFactoring.lean}

\lstinputlisting[language=Lean]{../11_CongruenceOfSquaresFactoring.lean}


\section{Quadruple Factor Theory}
\label{app:12_QuadrupleFactorTheory}
\texttt{12_QuadrupleFactorTheory.lean}

\lstinputlisting[language=Lean]{../12_QuadrupleFactorTheory.lean}


\section{GCD Cascade Factor Extraction}
\label{app:13_GCDCascadeFactorExtraction}
\texttt{13_GCDCascadeFactorExtraction.lean}

\lstinputlisting[language=Lean]{../13_GCDCascadeFactorExtraction.lean}


\section{Pythagorean Tree Factoring Core}
\label{app:14_PythagoreanTreeFactoringCore}
\texttt{14_PythagoreanTreeFactoringCore.lean}

\lstinputlisting[language=Lean]{../14_PythagoreanTreeFactoringCore.lean}


\section{Tropical Geometry Foundations}
\label{app:15_TropicalGeometryFoundations}
\texttt{15_TropicalGeometryFoundations.lean}

\lstinputlisting[language=Lean]{../15_TropicalGeometryFoundations.lean}


\section{Lorentz Group Structure}
\label{app:16_LorentzGroupStructure}
\texttt{16_LorentzGroupStructure.lean}

\lstinputlisting[language=Lean]{../16_LorentzGroupStructure.lean}



\end{document}
